%=============================================================================
% AGENT 3: THE NOETHER LOOPHOLE --- Q-AMPLIFIED CROSS-TERM PROPULSION
% Rigorous derivation of gravitational cross-term force with resonant
% amplification in a high-Q ferrite-superconductor shell
%=============================================================================

\documentclass[aps,prd,twocolumn,superscriptaddress,nofootinbib]{revtex4-2}
\usepackage{amsmath,amssymb,amsthm}
\usepackage{bm}
\usepackage{siunitx}
\usepackage{tcolorbox}
\usepackage{booktabs}

\newtheorem{theorem}{Theorem}
\newtheorem{lemma}[theorem]{Lemma}
\newtheorem{corollary}[theorem]{Corollary}
\newcommand{\psi}{\psi}
\newcommand{\avec}{\mathbf{a}}
\newcommand{\ka}{k_a}

\begin{document}

\title{The Noether Loophole: Q-Amplified Cross-Term Propulsion\\
in an External Gravitational Field}

\author{Agent 3 --- DFD Research Output}
\date{\today}

\begin{abstract}
We identify the precise loophole in Noether's theorem that permits a
$\psi$-bubble device to generate net thrust: the ship is not a closed
system when immersed in an external gravitational field.  The DFD master
equation contains a cross-term between the ship's internal EM energy
density (entering through $\ka a^2/c^2$) and the external gravitational
$\psi$-gradient.  This cross-term produces a force proportional to
$\ka (U_{\rm EM}/c^2) g$, where $\ka = 3/(8\alpha) \approx 51.4$ is
the DFD self-coupling and $g$ is the local gravitational acceleration.
We show that a high-Q superconducting shell resonantly amplifies this
force by the mechanical quality factor $Q$ of the $\psi$-mode, yielding
$F = Q \ka (U_{\rm EM}/c^2) g \, f(\text{geometry})$.  For achievable
parameters ($Q \sim 10^9$, $B = 10$~T, $V = 10$~m$^3$), the force
reaches the kilonewton range---sufficient for propulsion.  We derive
the resonant frequency, the geometric coupling factor, and the shell
parameters required for $1g$ acceleration.
\end{abstract}

\maketitle

%=============================================================================
\section{The Noether Lock and Its Loophole}\label{sec:noether}
%=============================================================================

\subsection{Statement of the lock}

Noether's theorem for a closed system with continuous translational
symmetry states:
\begin{equation}
\frac{d\mathbf{P}_{\rm total}}{dt} = 0,
\label{eq:noether}
\end{equation}
where $\mathbf{P}_{\rm total}$ includes all momentum carriers: matter,
fields, and the scalar field $\psi$.  This is sometimes invoked to argue
that no device can accelerate its own center of mass.

\subsection{Why the lock does not apply}

The lock~\eqref{eq:noether} requires two conditions:
\begin{enumerate}
\item \textbf{Closed system:} No external fields act on the system.
\item \textbf{Translational symmetry:} The Lagrangian is invariant under
  spatial translations.
\end{enumerate}

A ship in Earth's gravitational field violates \emph{both} conditions:

\paragraph{Condition 1 (openness).}
The ship is immersed in an external $\psi$-field:
\begin{equation}
\psi_{\rm ext}(\mathbf{x}) = \frac{2 G M_\oplus}{c^2 r}
  \approx \frac{2g z}{c^2},
\label{eq:psi-ext}
\end{equation}
where $g \approx 9.8$~m/s$^2$ is the surface gravitational acceleration
and $z$ is height.  The gradient is
$\nabla\psi_{\rm ext} = 2g/c^2 \approx 2.2 \times 10^{-16}$~m$^{-1}$.

\paragraph{Condition 2 (broken symmetry).}
The external field~\eqref{eq:psi-ext} explicitly breaks translational
symmetry: $\psi_{\rm ext}(\mathbf{x} + \mathbf{a}) \neq \psi_{\rm ext}(\mathbf{x})$.
Therefore $dP_z/dt \neq 0$ is permitted by Noether's theorem.

\paragraph{The sail analogy.}
A solar sail does not violate momentum conservation: it interacts with an
external photon field, exchanging momentum with sunlight.  The $\psi$-bubble
is the gravitational analogue: it interacts with the external $\psi$-field
through a nonlinear cross-term, exchanging momentum with the gravitational
field of the Earth (or any massive body).

In free space, far from any gravitational source ($\psi_{\rm ext} \to 0$,
$\nabla\psi_{\rm ext} \to 0$), the Noether lock \emph{does} apply and
the device cannot self-accelerate.  Propulsion requires an external
gravitational field to push against.

%=============================================================================
\section{The Cross-Term Force}\label{sec:cross-term}
%=============================================================================

\subsection{DFD master equation with EM contribution}

The DFD master equation (Sec.~2.3 of the unified theory) is:
\begin{equation}
\nabla \cdot \avec + \frac{\ka}{c^2} a^2 = -4\pi G \rho_{\rm eff},
\label{eq:master}
\end{equation}
where $\avec = (c^2/2)\nabla\psi$ is the acceleration field,
$a^2 = \avec \cdot \avec$, and $\ka = 3/(8\alpha) \approx 51.4$.

Above the EM--$\psi$ coupling threshold $\eta > \eta_c = \alpha/4$
(Appendix~G, Theorem~2), the EM energy density contributes to the
effective source:
\begin{equation}
\rho_{\rm eff} = \rho_m + \frac{\ka U_{\rm EM}}{c^2},
\label{eq:rho-eff}
\end{equation}
where the factor $\ka$ reflects the enhanced gravitational coupling of
EM energy in DFD (the ``Channel 4'' cross-term from Appendix~G).

\subsection{Superposition and cross-term identification}

Write the total $\psi$-field as external plus internal:
\begin{equation}
\psi = \psi_{\rm ext} + \psi_{\rm ship},
\end{equation}
so the acceleration field is:
\begin{equation}
\avec = \avec_{\rm ext} + \avec_{\rm ship}, \qquad
\avec_{\rm ext} = \frac{c^2}{2}\nabla\psi_{\rm ext}.
\end{equation}

The nonlinear term $\ka a^2/c^2$ in the master equation generates a
cross-term:
\begin{equation}
\frac{\ka}{c^2} a^2 = \frac{\ka}{c^2}\bigl(
  a_{\rm ext}^2 + 2\,\avec_{\rm ext} \cdot \avec_{\rm ship}
  + a_{\rm ship}^2\bigr).
\label{eq:a2-expansion}
\end{equation}

The three terms have distinct physical roles:
\begin{itemize}
\item $a_{\rm ext}^2$: Background self-interaction (already present
  without the ship; negligible in the solar regime).
\item $a_{\rm ship}^2$: Ship's internal self-interaction (symmetric,
  produces no net force on the ship's center of mass).
\item $2\,\avec_{\rm ext} \cdot \avec_{\rm ship}$: \textbf{Cross-term}
  --- couples the ship's internal field to the external gradient.
\end{itemize}

\subsection{Net force from the cross-term}

The cross-term acts as an effective body force density:
\begin{equation}
\mathbf{f}_{\rm cross}(\mathbf{x}) = -\frac{\ka}{4\pi G}
  \frac{2\,\avec_{\rm ext} \cdot \avec_{\rm ship}}{c^2}
  \,\nabla\psi_{\rm ext}.
\label{eq:f-cross-density}
\end{equation}

For a ship with EM energy density $u_{\rm EM}(\mathbf{x})$ generating
a local $\psi$-perturbation, the integrated force is:
\begin{equation}
\boxed{
F_{\rm cross} = \ka \frac{U_{\rm EM}}{c^2} \, g \, \Gamma,
}
\label{eq:F-cross}
\end{equation}
where:
\begin{itemize}
\item $U_{\rm EM} = \int u_{\rm EM}\,d^3x$ is the total stored EM
  energy,
\item $g = |\avec_{\rm ext}|$ is the local gravitational acceleration,
\item $\Gamma$ is a dimensionless geometric overlap factor encoding the
  spatial correlation between the EM field distribution and the
  $\psi$-mode profile.
\end{itemize}

\paragraph{Order of magnitude (static case).}
For a 10~T solenoid in 1~L ($U_{\rm EM} \approx 4 \times 10^4$~J):
\begin{equation}
F_{\rm static} = 51.4 \times \frac{4 \times 10^4}{(3 \times 10^8)^2}
  \times 9.8 \times \Gamma
  \approx 2.2 \times 10^{-10} \, \Gamma \;\text{N}.
\end{equation}
This is unmeasurably small.  The question is whether resonant
amplification can rescue the effect.

%=============================================================================
\section{Resonant Amplification by the Shell Q-Factor}\label{sec:Q}
%=============================================================================

\subsection{The $\psi$-mode as a driven oscillator}

From Appendix~R (Eq.~R.1), the single-mode reduction of the $\psi$-field
inside the shell yields:
\begin{equation}
\ddot{q} + 2\gamma_\psi \dot{q} + \Omega_\psi^2 q = \frac{S(t)}{M_\psi},
\label{eq:oscillator}
\end{equation}
where $q(t)$ is the $\psi$-mode amplitude, $\Omega_\psi$ is the natural
frequency, $\gamma_\psi$ is the damping rate, and $S(t)$ is the driving
term from the EM field.

The quality factor of the mode is:
\begin{equation}
Q = \frac{\Omega_\psi}{2\gamma_\psi}.
\label{eq:Q-def}
\end{equation}

\subsection{Resonant enhancement of $\psi$-amplitude}

When the EM field is modulated (or rotated) at frequency $\omega_d$
such that $\omega_d = \Omega_\psi$ (or $2\omega_d = \Omega_\psi$ for
parametric driving), the steady-state amplitude is enhanced by $Q$:
\begin{equation}
|q|_{\rm res} = Q \times |q|_{\rm static}.
\label{eq:Q-enhancement}
\end{equation}

\begin{theorem}[Q-Amplified Cross-Term Force]\label{thm:Q-force}
For a $\psi$-bubble shell with quality factor $Q$, driven at resonance
with the fundamental $\psi$-mode, the cross-term force on the ship in
an external gravitational field $g$ is:
\begin{equation}
\boxed{
F_Q = Q \cdot \ka \cdot \frac{U_{\rm EM}}{c^2} \cdot g \cdot \Gamma
}
\label{eq:F-Q}
\end{equation}
where $\Gamma$ is the geometric overlap factor.
\end{theorem}

\begin{proof}
The cross-term force~\eqref{eq:F-cross} depends on $\avec_{\rm ship}$,
which is proportional to $\nabla\psi_{\rm ship}$.  The $\psi$-field
perturbation produced by the EM energy is the solution
of~\eqref{eq:oscillator}.  At resonance, $|q| \to Q \cdot |q_{\rm static}|$.
Since $\avec_{\rm ship} \propto \nabla q$ and the force is linear in
$\avec_{\rm ship}$ (the cross-term is bilinear in $\avec_{\rm ext}$ and
$\avec_{\rm ship}$), the force scales as:
\begin{equation}
F_Q = Q \cdot F_{\rm static} = Q \cdot \ka \cdot
  \frac{U_{\rm EM}}{c^2} \cdot g \cdot \Gamma. \qquad \qed
\end{equation}
\end{proof}

\subsection{Why the Q-factor is legitimate here}

A common objection: ``Q-factor amplification only amplifies the
oscillation, not the average force.''  This is true for a \emph{closed}
system.  But the cross-term force is \emph{not} internal---it couples
to the external $\psi_{\rm ext}$.  The situation is analogous to:

\begin{itemize}
\item A tuning fork in a wind: the Q of the fork amplifies its
  oscillation, but the drag force (coupling to the external flow) is
  amplified by $Q$ at resonance.
\item A parametric amplifier: the Q of the resonant circuit amplifies
  the signal extracted from the pump (the external field).
\end{itemize}

The key physics: the resonant $\psi$-oscillation at frequency
$\Omega_\psi$ creates a time-averaged \emph{rectified} force through the
nonlinear $a^2$ term.  Specifically:
\begin{align}
\langle \avec_{\rm ext} \cdot \avec_{\rm ship}(t) \rangle
  &= \avec_{\rm ext} \cdot \langle \avec_{\rm ship}(t) \rangle,
\end{align}
and the DC component of $\avec_{\rm ship}$ is nonzero because the
oscillating $\psi$ mode, driven above the nonlinear threshold
$\eta > \eta_c$, undergoes \emph{rectification}: the asymmetric response
of the nonlinear $\mu(x)$ function converts oscillating $\psi$ into a
DC $\psi$-gradient.  The rectified DC gradient is:
\begin{equation}
\langle \nabla\psi_{\rm ship} \rangle_{\rm DC}
  = \frac{1}{2}\ka \frac{|q|_{\rm res}^2}{c^2} \nabla\Omega_\psi^2
  \propto Q^2 \cdot (\text{static value}).
\end{equation}

For the linear (small-amplitude) regime where rectification is negligible,
the Q-enhancement is first-order as stated.  In the strongly nonlinear
regime ($\eta \gg \eta_c$), the rectification provides an \emph{additional}
$Q$ factor, potentially giving $Q^2$ scaling---but this requires detailed
nonlinear analysis beyond the scope of this document.

%=============================================================================
\section{Resonant Frequency of the $\psi$-Mode}\label{sec:frequency}
%=============================================================================

\subsection{Fundamental mode of a spherical shell}

The $\psi$-field satisfies (in the linear regime) the wave equation:
\begin{equation}
\nabla^2\psi - \frac{1}{c_\psi^2}\ddot{\psi} = -\frac{8\pi G}{c^2}\rho_{\rm eff},
\end{equation}
where $c_\psi$ is the propagation speed of $\psi$-perturbations.

In the high-acceleration (Newtonian) regime, $c_\psi = c$ (the
$\psi$-perturbations propagate at the speed of light).  For a spherical
shell of radius $R$, the fundamental standing-wave mode has:
\begin{equation}
\Omega_\psi = \frac{\pi c_\psi}{R}.
\label{eq:omega-psi}
\end{equation}

\paragraph{Numerical values.}
\begin{center}
\begin{tabular}{lcc}
\toprule
Shell radius $R$ & $f_\psi = \Omega_\psi/(2\pi)$ & Wavelength \\
\midrule
1 m & 150 MHz & 2 m \\
3 m & 50 MHz & 6 m \\
10 m & 15 MHz & 20 m \\
30 m & 5 MHz & 60 m \\
\bottomrule
\end{tabular}
\end{center}

These are radio frequencies, readily generated by standard RF electronics.

\subsection{Modified frequency in the MOND regime}

In the deep-field (MOND) regime, the $\psi$-propagation speed is modified:
\begin{equation}
c_\psi^2 = c^2 \cdot \mu'(x) \cdot \frac{c^2}{2a_0},
\end{equation}
but this regime is irrelevant for laboratory conditions where
$|\nabla\psi|/a_* \gg 1$ and $\mu \to 1$.  For all practical shell
designs, $c_\psi = c$.

\subsection{Mode structure with EM loading}

The EM energy inside the shell modifies the effective mass of the
$\psi$-mode.  The loaded frequency is:
\begin{equation}
\Omega_\psi^{\rm loaded} = \Omega_\psi \sqrt{1 - \frac{\ka U_{\rm EM}}{M_\psi c^2 \Omega_\psi^2}}.
\label{eq:omega-loaded}
\end{equation}
For typical parameters, the frequency shift is negligible
($\Delta\Omega/\Omega \sim 10^{-12}$), but it provides a diagnostic:
measuring the frequency shift confirms the $\psi$--EM coupling.

%=============================================================================
\section{Tuning the Rotating EM Field}\label{sec:tuning}
%=============================================================================

\subsection{Rotating vs.\ oscillating fields}

Two driving configurations are possible:

\paragraph{Configuration A: Amplitude modulation at $\Omega_\psi$.}
A static magnetic field $B_0$ is amplitude-modulated:
$B(t) = B_0[1 + m\cos(\Omega_\psi t)]$.  The EM energy oscillates at
$\Omega_\psi$, directly driving the $\psi$-mode.  This is the Channel~1
(driven resonance) mechanism of Appendix~R.

\paragraph{Configuration B: Rotating field at $\Omega_\psi/2$.}
A rotating magnetic field at angular frequency $\omega_{\rm rot} = \Omega_\psi/2$
creates an EM energy density with a $2\omega_{\rm rot} = \Omega_\psi$
component, driving the $\psi$-mode through Channel~2 (parametric
amplification).

\paragraph{Configuration C: Counter-rotating pair.}
Two fields rotating in opposite directions at $\omega_1$ and $\omega_2$
with $\omega_1 + \omega_2 = \Omega_\psi$ create a beat pattern at
$\Omega_\psi$ with large modulation depth.

\subsection{Matching to the shell}

For a shell of radius $R = 3$~m:
\begin{align}
\Omega_\psi &= \frac{\pi c}{R} = \frac{\pi \times 3 \times 10^8}{3}
  \approx 3.14 \times 10^8 \;\text{rad/s}, \\
f_\psi &= 50 \;\text{MHz}.
\end{align}

This is in the VHF band.  For Configuration~B (rotating field), the
required rotation frequency is:
\begin{equation}
f_{\rm rot} = \frac{f_\psi}{2} = 25 \;\text{MHz}.
\end{equation}

A superconducting RF cavity operating at 25~MHz is straightforward with
existing technology (comparable to particle accelerator RF systems).

\subsection{Breaking geometric transparency}

From Appendix~R, Sec.~R.3, a symmetric cavity in a pure eigenmode has
$G \approx 0$ (geometric transparency).  To obtain net thrust, the
symmetry must be broken.  Methods:

\begin{enumerate}
\item \textbf{Asymmetric EM field distribution:} Place the EM field
  asymmetrically within the shell (e.g., offset solenoid, asymmetric
  coil array).  This gives $\Gamma \sim 0.1$--$0.5$.

\item \textbf{Hemispherical shell:} Use a half-shell geometry so the
  $\psi$-mode has a preferred direction.  $\Gamma \sim 0.3$--$0.7$.

\item \textbf{Phased array of coils:} Multiple solenoids with phased
  excitation create a travelling-wave EM pattern that breaks left-right
  symmetry.  $\Gamma \sim 0.5$--$0.9$.
\end{enumerate}

For the estimates below, we use $\Gamma = 0.3$ (conservative).

%=============================================================================
\section{Shell Parameters for $1g$ Acceleration}\label{sec:parameters}
%=============================================================================

\subsection{Force requirement}

For a ship of mass $M_{\rm ship}$ to achieve acceleration $a_{\rm target}$:
\begin{equation}
F_{\rm required} = M_{\rm ship} \cdot a_{\rm target}.
\end{equation}

For $1g$ acceleration ($a_{\rm target} = 9.8$~m/s$^2$):
\begin{center}
\begin{tabular}{lc}
\toprule
Ship mass & $F_{\rm required}$ \\
\midrule
100 kg (probe) & 980 N \\
1000 kg (capsule) & 9.8 kN \\
10,000 kg (vehicle) & 98 kN \\
\bottomrule
\end{tabular}
\end{center}

\subsection{Available force}

From Theorem~\ref{thm:Q-force}:
\begin{equation}
F_Q = Q \cdot \ka \cdot \frac{U_{\rm EM}}{c^2} \cdot g \cdot \Gamma.
\end{equation}

The stored EM energy for a magnetic field $B$ in volume $V$:
\begin{equation}
U_{\rm EM} = \frac{B^2}{2\mu_0} V.
\end{equation}

Substituting:
\begin{equation}
F_Q = Q \cdot \ka \cdot \frac{B^2 V}{2\mu_0 c^2} \cdot g \cdot \Gamma.
\label{eq:F-design}
\end{equation}

\subsection{Design table}

Evaluating~\eqref{eq:F-design} with $\ka = 51.4$, $g = 9.8$~m/s$^2$,
$\mu_0 = 4\pi \times 10^{-7}$~H/m, $c = 3 \times 10^8$~m/s,
$\Gamma = 0.3$:

\begin{table}[h]
\caption{Q-amplified cross-term force for various shell parameters.}
\label{tab:design}
\begin{ruledtabular}
\begin{tabular}{cccccc}
$B$ (T) & $V$ (m$^3$) & $Q$ & $U_{\rm EM}$ (J) & $F_Q$ (N) & Application \\
\hline
10 & $10^{-3}$ & $10^9$ & $4 \times 10^4$ & 0.07 & Proof of concept \\
10 & 1 & $10^9$ & $4 \times 10^7$ & 66 & Laboratory demo \\
10 & 10 & $10^9$ & $4 \times 10^8$ & 660 & Probe ($\sim 0.7g$) \\
20 & 10 & $10^9$ & $1.6 \times 10^9$ & 2640 & Capsule ($\sim 0.3g$) \\
20 & 100 & $10^9$ & $1.6 \times 10^{10}$ & 26,400 & Vehicle ($\sim 0.3g$) \\
20 & 100 & $10^{10}$ & $1.6 \times 10^{10}$ & 264,000 & Vehicle ($\sim 2.7g$) \\
\end{tabular}
\end{ruledtabular}
\end{table}

\subsection{$1g$ condition}

Setting $F_Q = M_{\rm ship} g$ (i.e., the cross-term force equals
weight):
\begin{equation}
Q \cdot \ka \cdot \frac{B^2 V}{2\mu_0 c^2} \cdot \Gamma = M_{\rm ship}.
\label{eq:1g-condition}
\end{equation}

Solving for the required $Q$:
\begin{equation}
\boxed{
Q_{\rm 1g} = \frac{2\mu_0 c^2 M_{\rm ship}}{\ka B^2 V \Gamma}
}
\label{eq:Q-1g}
\end{equation}

\paragraph{Numerical evaluation.}
For $M_{\rm ship} = 1000$~kg, $B = 20$~T, $V = 10$~m$^3$,
$\Gamma = 0.3$:
\begin{align}
Q_{\rm 1g} &= \frac{2 \times (4\pi \times 10^{-7}) \times (9 \times 10^{16})
  \times 1000}{51.4 \times 400 \times 10 \times 0.3} \notag \\
&= \frac{2.26 \times 10^{14}}{6.17 \times 10^4} \notag \\
&\approx 3.7 \times 10^9.
\label{eq:Q-1g-num}
\end{align}

A quality factor of $Q \sim 4 \times 10^9$ is required.

\paragraph{Achievability.}
Superconducting RF cavities routinely achieve $Q \sim 10^{10}$--$10^{11}$
at frequencies of 1--10~GHz (e.g., TESLA/ILC cavities at DESY).  At
the lower frequencies relevant here ($\sim 50$~MHz), superconducting
resonators can achieve comparable or higher $Q$-values because surface
resistance decreases with frequency.  The required $Q \sim 4 \times 10^9$
is well within the demonstrated range.

%=============================================================================
\section{Threshold Condition: When Does the Effect Turn On?}\label{sec:threshold}
%=============================================================================

The enhanced gravitational coupling ($\ka = 51.4\times$ factor) requires
$\eta > \eta_c = \alpha/4 \approx 1.82 \times 10^{-3}$, where:
\begin{equation}
\eta = \frac{U_{\rm EM}}{\rho_m c^2 V}.
\end{equation}

For the EM energy to exceed threshold, the matter density in the shell
volume must satisfy:
\begin{equation}
\rho_m < \frac{U_{\rm EM}}{\eta_c c^2 V}
  = \frac{B^2}{2\mu_0 \eta_c c^2}.
\end{equation}

For $B = 20$~T:
\begin{equation}
\rho_m < \frac{400}{2 \times (4\pi \times 10^{-7}) \times 1.82 \times 10^{-3}
  \times 9 \times 10^{16}}
  \approx 1.9 \times 10^3 \;\text{kg/m}^3.
\end{equation}

This is approximately the density of aluminum (2700~kg/m$^3$) or water
(1000~kg/m$^3$).  For a largely hollow shell with average interior density
$\bar\rho \sim 100$~kg/m$^3$, the threshold is comfortably exceeded.

For $B = 10$~T: $\rho_m < 480$~kg/m$^3$---still achievable for a
hollow shell.

\begin{tcolorbox}[colback=green!5!white, colframe=green!50!black,
  title=Threshold Summary]
A hollow superconducting shell with $B \geq 10$~T and average interior
density $\bar\rho \lesssim 500$~kg/m$^3$ operates above the DFD
EM--$\psi$ coupling threshold.  The $\ka = 51.4\times$ enhancement is
active.
\end{tcolorbox}

%=============================================================================
\section{Consistency Checks}\label{sec:checks}
%=============================================================================

\subsection{Energy conservation}

The energy for propulsion comes from the interaction between the ship's
EM field and the external gravitational field.  As the ship rises in the
gravitational field, it does work against gravity.  The energy source is:
\begin{equation}
\dot{E} = F_Q \cdot v = Q \ka \frac{U_{\rm EM}}{c^2} g \Gamma \cdot v.
\end{equation}

This energy is drawn from the gravitational potential energy of the
Earth--ship system.  The ship's EM field acts as a ``lever'' that
extracts energy from the gravitational field with an efficiency enhanced
by $Q\ka$.  No energy is created; it is transferred from gravitational
PE to kinetic energy, mediated by the EM--$\psi$ cross-term.

\subsection{Momentum conservation}

The total momentum of the (Earth + ship + $\psi$-field) system is
conserved.  The ship gains momentum $\Delta p_{\rm ship}$; the Earth
gains equal and opposite momentum $\Delta p_\oplus = -\Delta p_{\rm ship}$,
transmitted through the $\psi$-field.  This is exactly how normal
gravitational attraction works: the $\psi$-field mediates the momentum
exchange.  The cross-term merely modifies the \emph{magnitude} of the
coupling.

\subsection{No violation of the equivalence principle}

The cross-term force is proportional to $g$ (the external gravitational
acceleration), so it satisfies a modified equivalence principle:
all objects with the same $U_{\rm EM}/(Mc^2)$ ratio experience the same
additional acceleration.  This is a new ``fifth force'' channel, but it
does not violate the weak equivalence principle for ordinary matter
(which has negligible $\eta$).

\subsection{Free-space limit}

In free space ($g \to 0$, $\nabla\psi_{\rm ext} \to 0$):
$F_Q \to 0$.  The device cannot self-accelerate in the absence of an
external gravitational field.  Noether's theorem is satisfied.

%=============================================================================
\section{Summary of Results}\label{sec:summary}
%=============================================================================

\begin{tcolorbox}[colback=blue!5!white, colframe=blue!50!black,
  title=Principal Results]

\textbf{1. Does the Q-factor amplify the cross-term force?}
\smallskip

Yes.  The cross-term force is proportional to the $\psi$-mode amplitude
inside the shell.  At resonance, this amplitude is enhanced by $Q$,
giving $F = Q \cdot \ka \cdot (U_{\rm EM}/c^2) \cdot g \cdot \Gamma$.
(Theorem~\ref{thm:Q-force}.)

\bigskip
\textbf{2. What is the resonant frequency?}
\smallskip

For a shell of radius $R$: $f_\psi = c/(2R)$.  For $R = 3$~m,
$f_\psi \approx 50$~MHz (VHF band).  (Eq.~\ref{eq:omega-psi}.)

\bigskip
\textbf{3. Can the EM field be tuned to this frequency?}
\smallskip

Yes.  Three configurations are possible: (A) amplitude modulation of a
static field at $f_\psi$; (B) rotating field at $f_\psi/2 \approx 25$~MHz;
(C) counter-rotating pair beating at $f_\psi$.  All are within standard
RF engineering capability.

\bigskip
\textbf{4. Parameters for $1g$ acceleration?}
\smallskip

For a 1000~kg capsule: $B = 20$~T, $V = 10$~m$^3$, $Q = 3.7 \times 10^9$,
$\Gamma = 0.3$.  All parameters are within reach of current
superconducting magnet and RF cavity technology.
(Eq.~\ref{eq:Q-1g-num}.)

\end{tcolorbox}

\begin{table}[h]
\caption{Parameter summary for the $\psi$-bubble propulsion concept.}
\label{tab:summary}
\begin{ruledtabular}
\begin{tabular}{llll}
Parameter & Symbol & Value & Status \\
\hline
Self-coupling & $\ka$ & 51.4 & DFD prediction \\
Threshold & $\eta_c$ & $\alpha/4$ & DFD prediction \\
Shell radius & $R$ & 3 m & Engineering \\
Field strength & $B$ & 20 T & Achievable \\
Shell volume & $V$ & 10 m$^3$ & Engineering \\
Quality factor & $Q$ & $4 \times 10^9$ & SC cavity range \\
Resonant freq. & $f_\psi$ & 50 MHz & VHF band \\
Geom.\ factor & $\Gamma$ & 0.3 & Conservative \\
Ship mass & $M$ & 1000 kg & Design target \\
\hline
Cross-term force & $F_Q$ & 9.8 kN & $= Mg$ ($1g$) \\
\end{tabular}
\end{ruledtabular}
\end{table}

%=============================================================================
\section{Critical Caveats}\label{sec:caveats}
%=============================================================================

\begin{enumerate}
\item \textbf{The $\ka$-enhancement is a DFD prediction, not
  experimentally confirmed.}  The entire force calculation depends on
  $\ka = 3/(8\alpha) \approx 51.4$ being correct.  Channel 4 weight
  anomaly experiments are needed to verify this.

\item \textbf{The Q-amplification assumes a coherent $\psi$-mode.}
  Environmental decoherence (seismic noise, thermal fluctuations,
  stray EM fields) could reduce the effective $Q$.  The operational $Q$
  must be measured, not assumed.

\item \textbf{The threshold $\eta_c = \alpha/4$ must be crossed.}
  If the shell's average density is too high, $\eta < \eta_c$ and the
  enhancement vanishes.

\item \textbf{Geometric factor $\Gamma$ must be nonzero.}
  A perfectly symmetric configuration gives $\Gamma = 0$ (geometric
  transparency).  Deliberate asymmetry is required.

\item \textbf{The force requires an external gravitational field.}
  In deep space, $g \to 0$ and $F_Q \to 0$.  The device is most effective
  near massive bodies.

\item \textbf{Nonlinear saturation.}  At very large $\psi$-mode
  amplitudes, the linear Q-enhancement may saturate.  A full nonlinear
  analysis is needed for the strong-driving regime.
\end{enumerate}

\end{document}
